
\chapter{Токенизация PAN}\label{ch:tokenization}
Утечка базы сохранённых карт даёт злоумышленнику миллионы PAN.
\highlight{Украденный токен привязан к~конкретному ТСП, устройству или каналу
и~за~их пределами обычно бесполезен}.

\section{Назначение токенизации}\label{sec:token-problem}

PAN редко остаётся в~одной системе. Продавец хранит его для оплаты, PSP (Payment
Service Provider, платёжный провайдер)~--- для возвратов, подписочный
сервис~--- для~рекуррентных списаний. Каждая копия расширяет зону риска:
компрометация одной из~систем раскрывает PAN, прошедший через все остальные.

Heartland Payment Systems в~2008--2009 годах потеряла порядка 130~миллионов PAN
после того, как SQL-инъекция на~корпоративном веб-сервере открыла атакующим
доступ к~процессинговой сети, где они установили вредоносное ПО для~перехвата
карточных данных~\cite{doj-heartland-2009}. Утечка ударила не~по~одному
продавцу, а~по~процессинговому узлу, через который проходили карты множества
эмитентов и~клиентов. По~данным 10-K Heartland за~2010~финансовый год (подан
в~SEC 10~марта 2011) совокупные расходы и~урегулирования по~инциденту с~момента
раскрытия в~январе~2009 и~до~31~декабря 2010~года составили \$146,1~млн
до~страхового возмещения~\cite{heartland-sec-2013}. Каждый скомпрометированный
PAN означал потенциальный перевыпуск с~новым номером, новым пластиком
и~обновлением всех подписочных привязок у~держателя карты. Эмитент платит
за~перевыпуск напрямую, держатель карты~--- прерванными подписками и~ручным
обновлением реквизитов.

Если вместо PAN по~маршруту передаётся токен, утечка базы ТСП не~раскрывает настоящий
PAN. Украденный токен бесполезен вне~своего домена: он привязан к~конкретному
ТСП, устройству или каналу и~в~чужом контексте
не~сработает~\cite{emvco-payment-tokenisation}.

Токенизация также сужает периметр PCI~DSS (Payment Card Industry Data Security
Standard, подробно~--- в~главе~\ref{ch:pci-dss}), но~только при~разделении
доступа. Освобождение от~требований получают компоненты, которые работают
исключительно с~токенами и~не~имеют доступа к~детокенизации. Хранилище токенов
и~любая система с~доступом к~PAN остаются в~периметре.
Подробнее о~периметре CDE (Cardholder Data Environment, среда обработки
карточных данных), сужении периметра, P2PE (Point-to-Point Encryption)
и~разновидностях токенизации (сетевая, токен PSP, vault и~FPE,
Format-Preserving Encryption)~--- в~разделе~\ref{sec:pci-scope-reduction}.

\begin{figure}[tbp]
  \centering
  \includefiguresvg[width=\linewidth]{assets/figures/ch09-tokenization/token-risk-before-after}
  \caption{Узлы~$\times$~сценарии токенизации; цвет ячейки~--- PCI-периметр узла.}\label{fig:token-risk-before-after}
\end{figure}

Сетевая и~локальная токенизация различаются охватом. Локальный токен действует
внутри сервиса, который его выпустил и~может детокенизировать.
\highlight{Сетевой токен выпускает и~распознаёт сама карточная сеть, поэтому он
  проходит весь путь от~ТСП до~эмитента и~уменьшает поверхность компрометации
на~каждом узле маршрута}. Матрица на~рис.~\ref{fig:token-risk-before-after}
показывает узлы по~периметру~PCI: \textit{в~CDE}~--- полный набор требований
PCI~DSS, \textit{расширенный~CDE}~--- узел содержит vault детокенизации,
\textit{вне~CDE}~--- узел видит только токен или DPAN. Vault в~сетевой
токенизации существует внутри одной карточной схемы: VTS у~Visa, MDES
у~Mastercard, TSP~НСПК у~Мира. Единого хранилища на~все схемы нет: один и~тот
же PAN у~разных схем токенизируется независимо.

\section{Разночтения термина \enquote{токен}}\label{sec:token-terminology}

\highlight{Под словом \enquote{токен} разные источники понимают разные объекты,
  а~семейство \texttt{*PAN} у~EMVCo, PCI~SSC, Visa, Mastercard и~НСПК
использует несовместимые определения}. Без~выбранного словаря разбор
архитектуры превращается в~сверку понятий на~каждом шаге.

PCI~SSC определяет токенизацию широко~--- это любой суррогат исходного PAN:
vault, FPE, сетевой токен. EMVCo Payment Tokenisation Framework, напротив,
описывает только сетевой токен с~собственным BIN из~выделенного сетью диапазона
и~детокенизацией через TSP; vault и~FPE в~эту область
не~входят~\cite{emvco-payment-tokenisation,pci-tokenization-guidelines}. Visa
и~Mastercard в~публичных материалах используют свои бренды (Visa Token, MDES,
MDES for Merchants, Visa Token Management Service)~--- они конкретизируют
EMVCo, не~расширяют его.

В~правилах НСПК запись \enquote{Токен~(DPAN)} повторяется в~каждом упоминании:
\enquote{DPAN} там~--- второе имя \enquote{Токена}, а~не~подтип с~областью
действия устройства. Это противоречит EMVCo, где \texttt{token}~--- родовой класс,
а~DPAN и~MPAN~--- подтипы с~областью действия устройства или ТСП. \texttt{MPAN}
в~канонической документации НСПК как термин отсутствует, и~функционального
аналога нет: Mir~Pay~SDK app2app работает поверх DPAN (выпуск
  DPAN~(\textit{provisioning}) в~Mir~Pay из~приложения банка плюс
\textit{In-Application e-commerce}~--- оплата из~приложения ТСП), а~токена
с~областью действия ТСП под именем \texttt{MPAN} в~правилах ПС~Мир
не~определено~\cite{nspk-tsp}.

Терминология TSP воспроизводит ту~же асимметрию: у~Visa и~Mastercard это VTS
и~MDES, формально TSP по~EMV Payment Tokenisation Specification~--- Technical
Framework. В~документации НСПК слово \enquote{TSP} в~этом смысле
не~применяется~--- функционально аналогичную роль для~Mir~Pay играет
\enquote{Провайдер} (АО~НСПК).

\begin{table}[tbp]
  \begingroup
  \footnotesize
  \setlength{\tabcolsep}{3pt}
  \renewcommand{\arraystretch}{1.25}
  \begin{tabularx}{\textwidth}{
      @{}P{0.12\textwidth}
      Y Y Y Y Y Y@{}
    }
    \toprule
    \headrow
    \thd{Термин} & \thd{EMVCo} & \thd{PCI~SSC} & \thd{Visa} &
    \thd{Mastercard} & \thd{Речь} & \thd{НСПК} \\
    \midrule
    token       & сетевой только & любой суррогат & Visa Token & MDES token & любой суррогат & \enquote{Токен (DPAN)} \\
    DPAN        & область действия устройства & --- & Visa Token (device) & MDES (device) & = network token & синоним \enquote{Токена} \\
    MPAN        & область действия ТСП & --- & VTS / TMS & MDES for Merchants (M4M) & vault CoF & --- \\
    FPAN        & underlying PAN & --- & funding source & underlying card & реальный PAN & \enquote{Номер Карты (FPAN)} \\
    pseudo PAN  & --- & --- (канон~--- \enquote{surrogate value}) & --- & --- & нестрого = токен & --- \\
    surrogate   & --- & surrogate value (канон) & --- & --- & редко & --- \\
    TSP         & Token Service Provider & --- & VTS & MDES & плавает & \enquote{Провайдер} Mir~Pay \\
    \bottomrule
  \end{tabularx}%
  \endgroup
  \caption{Семейство \texttt{*PAN} по~источникам.}\label{tab:token-terminology}
\end{table}

Прочерки в~колонке НСПК~--- сигнал, что термина в~канонической российской
документации нет: \enquote{pseudo PAN} и~\enquote{MPAN} встречаются только у~PCI~SSC
и~западных сетей. Инженер, ищущий \enquote{MPAN у~Mir~Pay}, в~документации НСПК
ничего не~найдёт~--- релевантный объект описан под именами \enquote{Mir~Pay
In-Application} и~\enquote{app2app}.

Канон книги дальше по~тексту: \emph{токен}~--- сетевой токен EMVCo, если
не~сказано иное; \emph{токен в~vault} и~\emph{FPE} упоминаются явно с~уточнением
\enquote{суррогат для~хранения и~аналитики, не~для~маршрутизации};
\emph{TSP~НСПК} ссылается на~функциональную роль \enquote{Провайдера} Mir~Pay,
а~не~на~официальный российский термин.

С~10~марта 2022~года Visa и~Mastercard не~обслуживают карты российских банков,
а~их TSP-сервисы (VTS, MDES, M4M, TMS) российскому эмитенту и~эквайеру
недоступны: новые DPAN и~MPAN на~Visa/Mastercard российского выпуска
не~генерируются, существующие DPAN в~Apple Pay и~Google Pay прекратили работу
в~тот~же день. В~России с~этого дня единственный действующий сетевой
TSP~--- TSP~НСПК для~карт~\enquote{Мир}: Mir~Pay (DPAN, область действия
устройства) и~Mir~Pay~SDK app2app для~\textit{In-Application e-commerce}.
UnionPay в~РФ процессит НСПК, но~токенизация UnionPay принадлежит UnionPay
International и~через TSP~НСПК не~проходит. Дальнейший разбор VTS, MDES, M4M,
TMS используется как справочный канон EMVCo: российскому эмитенту он даёт
терминологию для~разговора с~международными вендорами, но~не~доступ к~этим
инструментам.

\section{Структура сетевого токена}\label{sec:token-anatomy}

Сетевой токен несёт больше данных, чем обычный PAN: к~номеру прикреплены
метаданные о~запросчике, проверке держателя карты и~разрешённом
домене~\cite{emvco-payment-tokenisation-guide}.

Первый атрибут~--- \textbf{Token Requestor}, участник, запросивший токен:
кошелёк, ТСП, PSP или подключённое устройство. В~модели EMVCo TRID (Token
Requestor ID) назначает Token Service Provider при~регистрации запросчика.
Первые три цифры идентификатора~--- TSP Code, который выдаёт сама EMVCo,
оставшиеся~--- порядковый номер запросчика у~конкретного TSP. По TRID и~схема,
и~эмитент определяют, через какой канал выпущен токен. Одна карта может
породить несколько токенов с~разными сценариями и~жизненными циклами: отдельно
для~телефона, для~браузера,
для~подписки~\cite{emvco-payment-tokenisation-guide}.

Второй набор метаданных связывает токен с~качеством проверки держателя карты.
В~терминологии EMVCo это \textbf{Token Assurance Method}~--- характеристика
того, как программа токенизации оценила проверку держателя карты при~выпуске.
Отдельные схемы затем отображают эту характеристику в~свои уровни и~коды. Чем
строже проверка на~выпуске, тем выше доверие к~токену в~дальнейших
операциях~\cite{emvco-payment-tokenisation-guide}.

\textbf{Token Domain Restriction Controls} задают контекст, в~котором токен
разрешено использовать: NFC, интернет-покупки у~конкретного ТСП или
рекуррентные списания у~определённого получателя. Попытку вывести токен
за~пределы домена сеть отклоняет до~этапа
авторизации~\cite{emvco-payment-tokenisation,emvco-payment-tokenisation-guide}.

Проверка выполняется на~стороне TSP до~детокенизации: сервер
сопоставляет канал запроса (contactless, CNP~--- Card-Not-Present, операции
  без~физического предъявления карты, см.~главу~\ref{ch:3ds}; recurring~---
рекуррентное списание), TRID инициатора и~MID (Merchant ID, идентификатор ТСП
в~эквайринге) ТСП с~атрибутами ограничения, записанными при~выпуске токена. При
несовпадении TSP возвращает отказ на~шаге маршрутизации, до~того как запрос
дойдёт до~эмитента. Domain restriction отсекает операции в~инфраструктуре
токенизации, не~в~авторизации.

\textit{Matcher} выполняет поэтапное сравнение. У~токена~--- набор
атрибутов (разрешённые каналы, TRID, опциональный список MID), у~запроса~---
фактические значения. Реальные TSP проверяют шире (assurance method, expiry,
криптограмма, статус токена), но структура решения та же.

У~токена есть собственные операции жизненного цикла: публичные материалы Visa
перечисляют activate, resume, suspend и~delete, а~НСПК в~описании Mir~HCE (Host
Card Emulation, см.~главу~\ref{ch:contactless}) SDK отдельно описывает управление
жизненным циклом токена. Скомпрометированный токен останавливают отдельно,
без~перевыпуска самой
карты~\cite{emvco-payment-tokenisation}.\footnote{Публичные материалы схем:~\cite{visa-vts,nspk-tsp}.}

На~листинге ниже~--- учебная реализация \textit{matcher}'а. Запрос сравнивается
с~атрибутами токена; при любом несовпадении возвращается отказ с~причиной,
и~авторизационный запрос не~отправляется эмитенту. Учебная версия проверяет
три атрибута~--- канал, TRID и~опциональный список MID; промышленный TSP
добавляет проверку статуса токена, криптограммы (TAVV~--- Token Authentication
  Verification Value,
  криптограмма сетевого токена для~e-commerce, разобрана
в~разделе~\ref{sec:token-detokenization}), срока действия и~уровня assurance,
но~разбирает их той~же поэтапной проверкой.

\codefile{Python}{samples/ch09-token-domain/token_domain.py}

\section{Семейство токенов по~области действия}\label{sec:token-scope}

\highlight{Сетевые токены различаются по~области разрешённого использования:
устройство или ТСП}. Один FPAN (см.~раздел~\ref{sec:card-data-pan-variants})
порождает столько токенов, сколько у~держателя карты точек оплаты,
и~компрометация одного из~них не~разрушает остальные.

DPAN действует в~пределах одного устройства: один FPAN порождает столько DPAN,
сколько у~держателя карты устройств~--- iPhone, Apple~Watch, Pixel,
Galaxy~Wearable. Скомпрометированный DPAN отзывается отдельно; FPAN и~остальные
DPAN продолжают работать~\cite{emvco-payment-tokenisation,visa-vts}. Физическое
место хранения DPAN на~устройстве (SE, TEE, HCE) и~различия между Apple Pay,
Samsung Pay, Google Wallet, Mir~Pay, Garmin~/~Fitbit Pay разобраны
в~разделе~\ref{sec:contactless-wallets-by-trust}.

MPAN действует в~пределах одного ТСП: ТСП получает токен при~первой авторизации
и~использует дальше для~\textit{card-on-file}. У~Mastercard этот класс
реализован как MDES for Merchants (M4M); у~Visa аналогичная функциональность
доступна через VTS и~Token Management Service (Visa Acceptance Solutions).
\highlight{MPAN нельзя путать с~токеном в~vault у~PSP}: такой токен действует
внутри хранилища провайдера и~не~передаётся в~сеть, MPAN~--- полноценный сетевой
объект с~собственным TRID и~Token Domain Restriction Controls. У~НСПК аналога
нет: Mir~Pay~SDK app2app закрывает сценарий привязки к~устройству поверх DPAN
(выпуск DPAN и~\textit{In-Application e-commerce}), а~токен с~областью
действия ТСП под именем MPAN в~правилах ПС~Мир
не~определён~\cite{nspk-tsp}.

Токен в~vault на~стороне PSP~--- параллельная семья, не~EMVCo. Шлюз хранит PAN
в~собственном защищённом хранилище и~возвращает ТСП непрозрачную ссылку. Эта
ссылка не~маршрутизируется в~сети, не~требует TSP-детокенизации и~не~несёт
сетевого контекста (\textit{assurance}, \textit{domain restriction}).
Архитектурное различие с~сетевыми токенами разобрано
в~разделе~\ref{sec:token-network-vs-psp}.

Pseudo PAN стоит особняком: синтаксически валиден как PAN (длина и~алгоритм
Луна совпадают), но маршрутизируется не~в~сеть, а~в~vault или ключевую
инфраструктуру. Платёжным объектом не~является и~к~авторизации непригоден;
используется для~сужения CDE~\cite{pci-tokenization-guidelines,nist-sp-800-38g}.

\begin{table}[tbp]
  \begingroup
  \footnotesize
  \setlength{\tabcolsep}{3pt}
  \renewcommand{\arraystretch}{1.3}
  \begin{tabularx}{\textwidth}{
      @{}P{0.16\textwidth}
      Y Y Y Y Y@{}
    }
    \toprule
    \headrow
    \thd{Параметр} & \thd{FPAN} & \thd{DPAN} & \thd{MPAN} &
    \thd{Vault PSP} & \thd{Pseudo PAN} \\
    \midrule
    Эмитент токена   & банк-эмитент & TSP сети & TSP сети & PSP & организация \\
    Область действия & глобальная & одно устройство & один ТСП & один vault & один vault/ключ \\
    Маршрутизация    & в~сети & в~сети & в~сети & не~в~сети & не~в~сети \\
    Детокенизация    & --- & TSP перед эмитентом & TSP перед эмитентом & PSP перед сетью & vault/FPE-ключ \\
    Хранение         & эмитент, vault сети & TEE/SE устройства, vault сети & vault сети, у~ТСП & vault PSP & vault или ключи \\
    Отзыв            & перевыпуск карты & отдельно по~DPAN & отдельно по~MPAN & отдельно по~записи & ротация ключей \\
    Где встречается  & везде & Apple/\allowbreak Google/\allowbreak Mir~Pay & крупные \textit{card-on-file}-ТСП & PSP-интеграции & сужение CDE \\
    \bottomrule
  \end{tabularx}%
  \endgroup
  \caption{Семейство \texttt{*PAN} по~области действия и~маршрутизации.}\label{tab:token-scope-comparison}
\end{table}

\section{VTS, MDES, TSP НСПК}\label{sec:token-providers}

У~каждой крупной карточной сети свой токен-сервис, и~архитектура у~всех почти
одинакова: TSP работает с~Token Requestors, при необходимости запрашивает
подтверждение эмитента на~выпуск токена и~встраивает этот токен в~платёжный
поток вплоть до~авторизации. Различия касаются правил управления, жизненного
цикла и~интеграционной
политики~\cite{emvco-payment-tokenisation,emvco-payment-tokenisation-guide}.

\subsection*{VTS (Visa Token Service)}

Token Requestor (например, мобильный кошелёк) подаёт запрос на~выпуск токена.
Провайдер обращается к~эмитенту, чтобы подтвердить личность держателя карты,
и~при необходимости запрашивает дополнительную проверку. После одобрения Visa
заменяет PAN сетевым токеном и~передаёт его цифровому платёжному сервису
для~использования онлайн и~через NFC. По документации Visa такой токен
ограничивают конкретным устройством, интернет-продавцом или числом покупок
до~истечения~\cite{visa-vts}.

Эмитент в~VTS работает с~тремя группами процессов. ID\&V (Identification \&
Verification, проверка личности держателя карты при~выпуске токена)
и~дополнительная проверка относятся к~этапу выпуска. Token Inquiry оперирует
данными об~устройстве и~риске. Token Lifecycle покрывает операции activate,
resume, suspend и~delete~\cite{visa-vts}.

\subsection*{MDES (Mastercard Digital Enablement Service)}

Mastercard описывает MDES как единую платформу для~эмитентов, поставщиков
кошельков, ТСП и~других участников, запрашивающих токены. В~документации MDES
Token Connect среди открытых token requestors перечислены кошельки на~мобильных
и~встраиваемых устройствах, торговые платформы
и~ТСП~\cite{mc-mdes,mc-mdes-token-connect}.

Token Connect выделяет MDES среди прочих TSP: открытый протокол
push-provisioning позволяет эмитенту инициировать выпуск токена в~сторонний
кошелёк без~интеграции один-на-один с~каждым кошельком. У~VTS аналог~--- Visa
Token Service Push Provisioning~--- работает по~отдельной схеме интеграции
с~каждым приёмником.

\subsection*{TSP НСПК}

Публичные материалы НСПК раскрывают модель токенизации Мир скупее, чем Visa или
Mastercard, но описывают сценарии, в~которых Mir~HCE~SDK вызывает сервисы НСПК
для~токенизации карт Мир. Держатель карты подаёт заявку на~токенизацию,
указывает метод ID\&V, получает одноразовые ключи для~транзакций и~управляет
жизненным циклом токена~\cite{nspk-tsp}.

Национальный маршрут и~приложение Мир разобраны в~главе~\ref{ch:mir}.

\begin{table}[tbp]
  \begingroup
  \footnotesize
  \setlength{\tabcolsep}{4pt}
  \renewcommand{\arraystretch}{1.2}
  \begin{minipage}{\textwidth}
    \begin{tabularx}{\textwidth}{
        @{}P{0.28\textwidth}
        Y@{}
      }
      \toprule
      \headrow
      \thd{Компонент TSP} & \thd{Назначение} \\
      \midrule
      Token Requestor Interface & Связь с~ТСП, кошельком или провайдером \\
      Token Vault & Хранит соответствие токена и~PAN \\
      Issuer Interface & Взаимодействие с~эмитентом при~выпуске токена \\
      \bottomrule
    \end{tabularx}
  \end{minipage}%
  \endgroup
  \caption{Архитектурные компоненты TSP~карточной сети.}\label{tab:token-tsp-comparison}
\end{table}

\section{Сетевой токен и~токен PSP}\label{sec:token-network-vs-psp}

\highlight{Сетевой токен рассчитан на~путь от~точки покупки через
  эквайера и~платёжную сеть до~авторизации у~эмитента; локальная токенизация
  на~стороне PSP остаётся внутри более узкого периметра ТСП и~его
PSP}~\cite{emvco-payment-tokenisation}.

Локальный токен PSP убирает PAN из~базы ТСП и~заменяет его непрозрачной ссылкой
на~хранилище провайдера, но сетевого токенового контекста при~этом не~создаёт.
EMVCo, напротив, описывает сетевой токен как идентификатор, который можно
ограничить ТСП, устройством или платёжным
сценарием~\cite{emvco-payment-tokenisation,emvco-payment-tokenisation-guide}.

Эмитент получает с~сетевым токеном дополнительный контекст: кто запросил токен,
как была проверена личность держателя карты и~в~каком домене токен разрешено
использовать. EMVCo отдельно отмечает возможность управлять и~заменять сетевые
токены для~ТСП, устройств или типов транзакций, часто без~участия держателя
карты~\cite{emvco-payment-tokenisation}.

Локальная токенизация проще в~интеграции: весь обмен замкнут внутри PSP.
Сетевая токенизация требует прямой работы с~токен-сервисом сети или подключения
провайдера, который скрывает эту интеграцию от~ТСП.

Хранилищная (\textit{vault}) токенизация и~FPE для~защиты данных в~покое решают
другую задачу: заменяют PAN суррогатом в~базах аналитики, CRM или биллинга,
чтобы эти системы вышли из~периметра CDE. Карточная сеть здесь не~участвует,
и~такой суррогат непригоден для~транзакции.

Табл.~\ref{tab:token-network-vs-psp} сводит архитектурное различие к~практическим
последствиям для~интегратора: портативности токена, периметру PCI~DSS у~ТСП
и~PSP, жизненному циклу при~перевыпуске карты и~стоимости интеграции.

\begin{table}[H]
  \begingroup
  \footnotesize
  \setlength{\tabcolsep}{4pt}
  \renewcommand{\arraystretch}{1.2}
  \begin{tabularx}{\textwidth}{
      @{}P{0.23\textwidth}
      Y
      Y@{}
    }
    \toprule
    \headrow
    \thd{Параметр} &
    \thd{Сетевой токен (VTS / MDES / TSP НСПК)} &
    \thd{Токен PSP} \\
    \midrule
    Портативность & да & нет \\
    Область PCI~DSS ТСП & минимальная & сужена \\
    Область PCI~DSS PSP & сужена & полная \\
    Доля одобрений & по~маркетинговым материалам схем~--- выше\textsuperscript{*} & без~изменений \\
    Жизненный цикл & автообновление & ручное + Account Updater\textsuperscript{\dag} \\
    Стоимость интеграции & средняя & низкая \\
    \bottomrule
  \end{tabularx}%
  \endgroup
  \par\vspace{0.3em}
  \begin{minipage}{\textwidth}
    \footnotesize
    \textsuperscript{*}~По маркетинговым материалам Visa и~Mastercard;
    конкретное значение зависит от~сегмента ТСП и~не~имеет публичного
    независимого подтверждения.\par
    \textsuperscript{\dag}~Account Updater~--- сервис карточной сети,
    обновляющий PAN и~expiry у~ТСП при~перевыпуске карты.
  \end{minipage}
  \caption{Сетевой токен и~токен PSP по~шести параметрам.}\label{tab:token-network-vs-psp}
\end{table}

\section{Внутренняя токенизация банка}\label{sec:token-internal-bank}

\highlight{Внутренний vault нужен банку, когда PAN надо вынести из~собственных
  нижестоящих (\textit{downstream}) систем~--- CRM, аналитики, биллинга,
  антифрод-витрин,~--- но при~этом сохранить обратную операцию для~узкого круга
авторизованных сервисов}. Сетевой токен решает другую задачу~--- убирает PAN
из~авторизационного и~клирингового потока до~эмитента; внешний vault PSP
переносит хранение PAN за~периметр организации. Три механизма
не~взаимозаменяемы: банк, который пытается сократить периметр PCI~DSS одним
из~них, не~получит ожидаемого результата.

\subsection*{Vault со~справочной таблицей и~FPE}

\textbf{Vault со~справочной таблицей (\textit{lookup})} генерирует случайный
токен и~хранит соответствие PAN~$\leftrightarrow$~токен в~отдельном защищённом
хранилище. Безопасность опирается на~хранилище и~разграничение доступа: знание
токена без~доступа к~vault не~даёт PAN. PCI~SSC Tokenization Product Security
Guidelines (2015) выделяет Card Data Vault в~отдельный домен требований
(Domain~3): обязательное шифрование PAN внутри хранилища, управление доступом,
отдельный жизненный цикл ключей шифрования
по~ISO~11568-1~\cite{pci-tokenization-guidelines}.

\textbf{Format-Preserving Encryption (FPE)} по~NIST~SP~800-38G~--- другой
класс: PAN шифруется в~значение той~же длины и~алфавитного состава
без~отдельного хранилища соответствий, обратная операция~--- расшифровка тем~же
ключом. FF1~--- единственный режим FPE, сохранившийся в~текущей редакции NIST:
второй публичный черновик SP~800-38G Revision~1 (3~февраля 2025) полностью
исключил FF3 после работы Beyne о~слабости \textit{tweak schedule}, поднял
минимальный размер домена с~ста до~миллиона возможных входов, запретил обратную
функцию AES и~\textit{floating-point arithmetic} в~реализациях FF1
(последнее~--- из-за исторического бага Bouncy Castle, обнаруженного
Bleichenbacher)~\cite{nist-sp-800-38g}.

Vault и~FPE расходятся по~модели компрометации. В~vault компрометация хранилища
или ключа шифрования раскрывает все хранящиеся PAN, но~компрометация отдельной
записи ограничивается этой записью. \highlight{В~FPE компрометация
ключа~--- это компрометация всех когда-либо выданных токенов сразу}: каждый
из~них криптографически выводится из~PAN и~общего ключа, один секрет связывает
их все.

\subsection*{Архитектура vault внутри банка}

\textbf{Хранилище соответствий} хранит пары PAN~$\leftrightarrow$~токен. PAN
зашифрован: у~западного банка~--- AES-256-GCM на~ключах, управляемых KMS или
HSM; у~российской финорганизации в~CDE по~Положению Банка России
от~30.01.2025 №~851-П и~ГОСТ~Р~57580.1-2017~--- Кузнечик через
сертифицированный СКЗИ~\cite{cbr-851p,gost-57580-1}. PCI~SSC требует, чтобы
хранилище соответствий и~настоящие PAN находились в~PCI~DSS-периметре, даже
если все нижестоящие системы видят только
токены~\cite{pci-tokenization-guidelines}.

\textbf{HSM или KMS поверх HSM} управляет ключами шифрования хранилища
и~ключами генерации токенов, если те~используются. Для~физических HSM PCI~SSC
требует валидации по~FIPS~140-2 Level~3 и~работы в~FIPS-режиме;
для~\textit{cloud HSM-as-a-Service} действуют отдельные требования по~изоляции
между арендаторами и~передаче ключей. У~российской финорганизации HSM
сертифицирован ФСБ как СКЗИ класса КС2 или КС3; ключи шифрования хранилища
выводятся из~ключевой иерархии, разобранной в~гл.~\ref{ch:crypto}.

\textbf{API \texttt{tokenize} и~\texttt{detokenize}} разделены по~правам.
Системы, записывающие данные, имеют только право \texttt{tokenize}: передают
PAN, получают токен. Системы, выполняющие авторизацию или возврат, имеют право
\texttt{detokenize}: передают токен, получают PAN. Большинство нижестоящих
узлов не~получают ни~того, ни другого~--- они работают только с~уже выданными
токенами. Через это разделение CDE сужается до~vault и~узлов детокенизации;
остальные системы выходят из~периметра PCI~DSS.

\textbf{Аудит} хранит каждый вызов \texttt{tokenize} и~\texttt{detokenize}: кто
инициировал, какой токен затронут, для~\texttt{detokenize}~--- обязательная
причина. По~ГОСТ~Р~57580.1-2017 (направление РЗИ~раздел~8) финансовая
организация ведёт учёт объектов и~ресурсов доступа на~всех уровнях
информационной инфраструктуры~\cite{gost-57580-1}; по~Положению Банка России
от~30.01.2025 №~851-П~--- регистрирует действия работников и~клиентов
с~автоматизированными системами по~детализированному составу полей (дата,
время, идентификатор работника, код участка, результат)~\cite{cbr-851p}.

Облачный vault (AWS Payment Cryptography, Azure Confidential Computing, Google
Cloud HSM) разворачивается за~недели, но~для российского банка, которому надо
выполнять Положение Банка России от~30.01.2025 №~851-П и~ГОСТ~Р~57580.1-2017,
эта модель неприемлема: ключи и~хранилище оказываются за~пределами суверенного
периметра.
\textit{On-premises} vault с~собственным HSM~--- единственный рабочий вариант
для~российского банка; срок развёртывания на~порядок больше, сертификация ФСБ
и~приёмка требуют отдельного проекта.

\subsection*{Российская специфика: vault как единственный путь}

В~России FPE недоступен на~уровне сертифицированной криптографии. ТК~26
(Технический комитет по~стандартизации \enquote{Криптографическая защита
информации} при~Росстандарте) не~стандартизировал ни~один режим FPE~---
ни~на~Кузнечике (ГОСТ~34.12-2018), ни~на~AES. В~опубликованных Рекомендациях
по~стандартизации (серия Р~1323565.1) и~Методических рекомендациях ТК~26
на~2024--2026~годы FPE не~упоминается; тематически ближайший документ~--- МР
по~режиму работы блочных шифров для~защиты носителей информации
с~блочно-ориентированной структурой, который решает другую задачу (защита
данных на~диске, не~FPE короткого поля)~\cite{tc26-recommendations}.
Сертификата ФСБ на~FPE-конструкцию, построенную поверх Кузнечика или AES, нет
ни~у~одного российского производителя.

Следствие для~российского банка: внутренняя токенизация~--- это vault,
а~не~FPE. Конструкция: хранилище соответствий с~PAN, зашифрованным Кузнечиком
через сертифицированный СКЗИ; ГОСТ~TLS до~vault при~обращении по~сети; HSM или
СКЗИ в~доверенном сегменте сети; меры по~Положению Банка России от~30.01.2025
№~851-П и~ГОСТ~Р~57580.1-2017 в~зависимости от~уровня защиты, требуемого
регулятором для~конкретного банка~--- усиленный для~системно значимых
организаций и~операторов услуг платёжной инфраструктуры значимых платёжных
систем, стандартный
для~остальных~\cite{cbr-851p,gost-57580-1}.

\subsection*{Когда внутреннего vault достаточно, когда нет}

Внутренний vault сокращает периметр своих нижестоящих систем, но~не~выводит PAN
из~авторизационного и~клирингового потока до~эмитента~--- для~этого нужен
сетевой токен; и~не~переносит хранение PAN за~организацию~--- для~этого нужен
внешний PSP-vault.

\begin{table}[tbp]
  \begingroup
  \footnotesize
  \setlength{\tabcolsep}{4pt}
  \renewcommand{\arraystretch}{1.3}
  \begin{tabularx}{\textwidth}{
      @{}P{0.30\textwidth}
      Y Y Y@{}
    }
    \toprule
    \headrow
    \thd{Вопрос} & \thd{Внутренний vault} & \thd{Сетевой токен (DPAN, MPAN)} & \thd{Внешний PSP-vault} \\
    \midrule
    PAN не~должен передаваться до~эмитента в~открытом виде? & нет, vault подставляет PAN перед~авторизацией & да, запрос передаётся с~токеном & нет, PSP подставляет PAN перед~сетью \\
    Хранение PAN должно быть вне~организации? & нет, vault внутри банка & нет, vault карточной сети & да, у~PSP \\
    Сокращаем CDE нижестоящих систем (CRM, аналитика, биллинг)? & да, основная задача & частично, если эти системы видели токен, а~не~PAN & да, если ТСП передаёт PAN только PSP \\
    Контроль ключей по~851-П от~30.01.2025 и~ГОСТ~Р~57580.1-2017 на~стороне банка? & да & нет (у~TSP сети) & нет (у~PSP) \\
    Автообновление при~перевыпуске карты? & нет, нужен внешний Account Updater & да, через обновление жизненного цикла токена в~TSP & зависит от~PSP \\
    \bottomrule
  \end{tabularx}%
  \endgroup
  \caption{Внутренний vault, сетевой токен и~внешний PSP-vault: критерии выбора.}\label{tab:token-tokenization-choice}
\end{table}

Архитектура крупного российского банка использует все три механизма
одновременно: внутренний vault для~аналитики, CRM и~биллинга; сетевые токены
через TSP~НСПК для~карт~\enquote{Мир} (другие сетевые токены российскому
  эмитенту с~марта~2022 недоступны~--- VTS и~MDES не~обслуживают карты
  российских банков, а~UnionPay токенизирует на~стороне UnionPay International,
не~через НСПК); внешние PSP-vault для~\textit{e-commerce}-ТСП, где банк
выступает эквайером. Эти механизмы не~взаимозаменяемы: у~каждого свой
жизненный цикл, свой набор требований PCI~DSS и~своя точка ответственности
при~инциденте.

\section{Антипаттерн: BIN и~хвост вместо токена}\label{sec:token-bin-last4-antipattern}

\highlight{Локальная схема, где токен строится из~сохранённого BIN,
  случайных средних цифр и~последних четырёх цифр исходного PAN,
  остаётся маской PAN: CDE у~ТСП не~сокращается, а~утечки
коррелируются без~PAN и~без~доступа к~vault}. Такую маску выдают
за~локальную токенизацию у~небольших ТСП и~молодых PSP, потому что она
сохраняет привычное отображение на~чеке и~даёт маршрутизацию
по~карточной сети без~обращения к~vault.

\subsection*{Сколько PAN проходит проверку Луна под одной маской}

При 16-значном PAN со~стандартной маской \enquote{BIN~6 + хвост~4}
средняя часть содержит 6 цифр. Сумма Луна по~модулю~10~---
линейное уравнение от~цифр, удвоение в~алгоритме
Луна~--- биекция на~$\{0,\dots,9\}$ (0$\to$0, 1$\to$2, 2$\to$4,
3$\to$6, 4$\to$8, 5$\to$1, 6$\to$3, 7$\to$5, 8$\to$7, 9$\to$9). Для
любой комбинации $d$ свободных цифр сумма пробегает все десять
остатков равномерно; из~$10^d$ комбинаций ровно $10^{d-1}$ проходят
проверку Луна. Для маски 6+4 это $10^5 = 100\,000$ PAN
на~один публичный слепок (\textit{fingerprint}). Полная длина PAN
создаёт видимость пространства $10^{16}$; маска оставляет пять
десятичных порядков.

После ревизии ISO~7812 2017~года идентификатор эмитента занимает
восемь цифр; требование 3.4.1 PCI~DSS ограничивает отображение PAN
полным BIN и~последними четырьмя цифрами, а~FAQ~1492 разъясняет это
правило для~8-значного BIN~\cite{iso-7812-transition}.\footnote{Требование 3.4.1 PCI~DSS и~FAQ~1492 о~8-значном BIN:~\cite{pci-dss-4,pci-faq-1492}.}
У~стандартного 16-значного PAN свободная часть сжалась до~четырёх цифр,
и~количество кандидатов на~маску упало до~$10^3 = 1\,000$. Расширение
BIN увеличивает адресное пространство эмитентов и~одновременно сужает
множество кандидатов у~маски, выданной за~токен, в~сто раз.

Для любой длины PAN ответ равен $10^{d-1}$, где $d$~--- число свободных
цифр в~средней части (для шестизначного BIN и~четырёхзначного хвоста:
  13-значный устаревший PAN Visa~--- $10^2$, 15-значный Amex~--- $10^4$,
19-значный длинный PAN~--- $10^8$).
Формула проверяется прямым перебором: при~учебной длине 8 цифр
с~маской 3+2 формула даёт $10^2$, и~перебор всех $10^3$ комбинаций
средней части даёт то~же число.

Код ниже реализует ровно этот антипаттерн: \texttt{tokenize\_by\_bin\_tail}
получает длину BIN и~хвоста, генерирует свободные средние цифры
и~подбирает последнюю свободную цифру так, чтобы суррогат проходил проверку
Луна. Vault намеренно не~моделируется: даже если связь суррогат~$\to$~FPAN
хранится отдельно, публичное значение уже сохраняет слепок исходного PAN.

\codefile{Python}{samples/ch09-bin-last4/bin_last4.py}

\subsection*{Корреляция утечек без~знания PAN}

Маска \texttt{(BIN, last4)}~--- детерминированная функция PAN без~ключа.
Один и~тот~же PAN у~двух разных ТСП даёт одинаковые BIN и~хвост;
две независимые утечки баз токенов сшиваются по~равенству
маски без~знания исходного PAN, без~доступа к~vault и~без~криптоанализа.
\highlight{Свойство \textit{unlinkability}, ради которого токенизация
  существует как механизм защиты PAN, у~такой схемы отсутствует
на~уровне определения}~\cite{pci-tokenization-guidelines}. Вероятность
случайного совпадения BIN и~хвоста у~двух независимых PAN равна
$10^{-10}$ в~равномерной модели.\footnote{Равномерная модель трактует все
  шестнадцать цифр PAN как независимые равномерные на~$\{0,\ldots,9\}$.
  Контрольная цифра Луна определена остальными пятнадцатью, поэтому last4
  содержит три независимые цифры, а~не~четыре; на~оценку $10^{-10}$ это
  не~влияет, так как контрольная сумма распределена равномерно
  по~своим префиксам. Реальные BIN неравномерно распределены по~$10^6$
  комбинациям~--- сконцентрированы на~действующих эмитентах,~--- и~фактическая
  вероятность совпадения BIN на~несколько порядков выше: $10^{-10}$~---
нижняя оценка.} Для одной известной маски и~корпуса
из~$10^7$ карт ожидаемое число ложных совпадений равно $10^{-3}$;
массовую склейку утечек ограничивает распределение BIN, потому что
секретности в~схеме уже нет.

\subsection*{Метаданные транзакций сокращают поиск}

Yves-Alexandre de Montjoye и~соавторы показали на~трёхмесячном корпусе
из~1{,}1~миллиона держателей карт, что четыре пары
\emph{(дата, торговая точка)} однозначно идентифицируют 90\,\%
участников~\cite{de-montjoye-credit-2015}. В~их корпусе нет ни~PAN,
ни~BIN, ни~хвоста. Когда утечка добавляет к~метаданным слепок
\texttt{(BIN, last4)}, атакующий начинает с~малой группы кандидатов
по~самой маске; дата и~торговая точка нужны уже для~редких совпадений
после первичного поиска.

\subsection*{Допустимая роль маски}

Сохранённые BIN и~хвост~--- это \emph{маска для~отображения} (чек,
письмо, интерфейс, уведомление в~приложении), разрешённая PCI~DSS~4.0.1
при~условии, что больше нигде вне~CDE PAN
не~хранится~\cite{pci-dss-4,pci-faq-1492}.
Как \emph{суррогат для~хранения} такая конструкция не~работает:
по~PCI~DSS Tokenization Guidelines суррогат выводит систему из~CDE
только если необратим без~доступа к~vault и~если сам vault
изолирован~\cite{pci-tokenization-guidelines}; маска
\texttt{(BIN, last4)} не~удовлетворяет ни~одному из~двух условий.

BIN и~хвост в~суррогате нужны для~маршрутизации по~карточной сети
и~диапазону карточного продукта без~обращения к~vault. Эту задачу
решает справочная таблица BIN рядом с~непрозрачным токеном из~vault;
сам суррогат остаётся случайным. Выбор зависит от~назначения:
\begin{itemize}
  \item для~отображения~--- та~же маска BIN и~хвост, но~без~хранения
    средней части (короткая строка \enquote{BIN, четыре звёздочки,
    last4}, не~суррогат);
  \item для~внутреннего vault~--- случайный непрозрачный токен
    без~привязки к~цифрам PAN, либо FPE FF1 на~криптографическом
    ключе~\cite{nist-sp-800-38g};
  \item для~авторизационного потока~--- сетевой токен (DPAN или MPAN)
    через TSP схемы~\cite{emvco-payment-tokenisation}.
\end{itemize}

\section{Детокенизация и~доступ к~PAN}\label{sec:token-detokenization}

Эмитенту нужен настоящий PAN для~авторизации. В~модели EMVCo сетевой токен
проходит четыре звена от~точки покупки
до~эмитента~\cite{emvco-payment-tokenisation}:

\begin{enumerate}
  \item \textbf{Продавец} отправляет транзакцию с~сетевым токеном и~не~обязан
    хранить PAN карты;
  \item \textbf{Эквайер / PSP} передаёт запрос дальше как обычную платёжную
    операцию, но уже в~токенизированном виде;
  \item \textbf{Токен-сервис сети / TSP} распознаёт токен, применяет
    ограничения токена и~готовит запрос на~авторизацию по~PAN;
  \item \textbf{Эмитент} получает запрос на~авторизацию по~PAN вместе
    с~контекстом токена и~может учесть, через какой сервис прошёл выпуск
    и~насколько надёжно был подтверждён держатель карты.
\end{enumerate}

Продавец передаёт эквайеру токен (DPAN, Device PAN,
см.~раздел~\ref{sec:contactless-mobile}), срок действия токена, криптограмму
(TAVV, Token Authentication Verification Value, криптограмма сетевого токена
для~e-commerce; бесконтактную криптограмму для~NFC) и~свой MID. Эквайер
формирует авторизационный запрос (MTI~0100) и~направляет его в~карточную сеть.
Сеть распознаёт диапазон BIN (Bank Identification Number, банковский
идентификационный диапазон PAN) как принадлежащий токенам и~маршрутизирует
запрос в~TSP. TSP проверяет статус токена (active, не~suspended и~не~deleted),
domain restriction (канал запроса совпадает с~разрешённым доменом, TRID и~MID
соответствуют записи выпуска) и~валидность криптограммы. При успехе TSP
извлекает из~vault исходный FPAN (Funding PAN, исходный PAN, к~которому
привязан токен) и~expiry, присоединяет Token Assurance Level и~направляет
детокенизированный запрос эмитенту. Эмитент видит FPAN и~дополнительный
контекст: через какой Token Requestor пришла операция и~с~каким уровнем
проверки держатель карты был подтверждён при~выпуске токена. Ответ эмитента
возвращается через TSP, который заменяет FPAN обратно на~токен, прежде чем
передать ответ эквайеру, так что продавец не~видит настоящий PAN. Значения
на~рис.~\ref{fig:token-detokenization-flow} учебные: реальные Token~BIN range
у~VTS, MDES и~TSP~НСПК~--- закрытая часть документации.

\begin{figure}[t]
  \centering
  \includefiguresvg[width=\linewidth]{assets/figures/ch09-tokenization/token-detokenization-flow}
  \caption{Токенизированная авторизация: подмена DPAN~$\leftrightarrow$~FPAN на~стороне TSP.}\label{fig:token-detokenization-flow}
\end{figure}

PAN восстанавливается там, где он нужен для~авторизационного решения. Если токен
приостановлен, удалён или вышел из~допустимого состояния, транзакция не~доходит
до~авторизации: для~карты это прямое следствие состояния
токена~\cite{emvco-payment-tokenisation,visa-vts}.

Публичные документы EMVCo стандартизуют именно общую модель авторизации,
а~детали клиринга и~представления данных о~токене зависят от~правил конкретной
сети, поэтому из~общей схемы EMVCo их вывести нельзя.

Первый сценарий: TSP недоступен или отвечает с~таймаутом. Карточная сеть
не~может детокенизировать запрос и~возвращает эквайеру системный отказ.
Каскадный повтор через другого эквайера не~помогает: точка отказа~--- общий TSP
схемы, а~не~маршрут эквайера. ТСП может повторить запрос через короткий
интервал или, если FPAN доступен, отправить транзакцию без~токена ценой
расширения периметра PCI~DSS у~ТСП.

Второй сценарий: токен приостановлен или удалён между моментом формирования
запроса и~его обработкой в~TSP. Это возможно при~блокировке карты, когда
эмитент отправляет suspend, а~транзакция в~процессе обработки
(\textit{in-flight}) прошла через эквайера. TSP проверяет статус, видит
suspended и~отклоняет операцию: продавец получает отказ, хотя секундой ранее
тот же токен работал.

Третий сценарий: подписочный MIT (merchant-initiated transaction). Карта
перевыпущена, PAN изменился, но~сетевой токен автоматически переключен на~новый
FPAN через обновление жизненного цикла. Если обновление ещё не~завершилось
к~моменту очередного списания, TSP детокенизирует в~старый FPAN, по~которому
эмитент уже выдаёт отказ из-за закрытой карты. После синхронизации проблема
исчезает, но~первое списание после перевыпуска может потребовать повтор через
несколько часов.

Сквозная работа сетевого токена через нескольких эквайеров требует, чтобы все
участники цепочки поддерживали поля токена и~его ограничения по~домену; иначе
токен на~одном из~переходов преобразуется обратно в~PAN или операция
отклоняется.

Во~всех трёх сценариях запрос с~сетевым токеном останавливается на~стороне TSP,
и~эквайер не~получает PAN. Практический вывод: продавец после системного отказа
TSP может повторить операцию по~PAN, но~рискует двойным списанием после
восстановления сервиса. В~ISO~8583 идемпотентность повторов строится
на~комбинации RRN~(DE~37) и~STAN~(DE~11) либо на~Transaction Life Cycle
Identifier из~редакции~2003, а~в~PSP-API передаётся отдельный
\textit{idempotency key}, генерируемый клиентом.

Эквайер выделяет TSP-отказы (timeout, token suspended, domain mismatch)
отдельной метрикой. Формально такие коды возвращаются в~том~же DE~39, что
и~отказы эмитента, но~приходят со~специальными значениями (у~Visa~--- 5C, 94
и~аналогичные; плюс дополнительные поля сети, DE~44 и~DE~62), поэтому отличать
их от~эмитентских надо по~значению, а~не~по~полю: DE~39 у~них тот~же.

Эмитент сводит долю операций, не~дошедших до~авторизации из-за состояния
токена, в~ту~же витрину, что и~отказы по~стоп-листу: для~держателя карты
результат один~--- карта числится рабочей, операция не~прошла,~--- а~корневая
причина различна. Повторы после отказа разобраны
в~разделе~\ref{sec:decline-retry}.

Связь токенизации и~3-D~Secure разобрана в~разделе~\ref{sec:3ds-tokenization},
влияние токенизации на~периметр PCI~DSS~---
в~разделе~\ref{sec:pci-scope-reduction}, на~антифрод~---
в~главе~\ref{ch:fraud}.
